\vpage{1}{Three ways to reduce a model}{PAPER VI / A WORKED VISUAL ATLAS}{Chapter VI, Part I: Sections 1--3, 8, and 13. Read the three operations separately.}
This guide follows small objects that can be computed completely. A four-state chain connects packet transfer, Markov aggregation, and Bayesian conditioning. A two-state coupled evolution exposes memory. A three-node Dirichlet problem makes the Green-operator normalization visible. These models are related by the chapter's language, not by pretending that their state spaces are identical.
\vfig{three_views}
\captiontext{Left and center use one probability vector and its two block masses. The right panel is a separate retained/hidden model, developed on \pagerefV{22}. Every curve is drawn from an explicit formula.}
\head{Coordinates retain; aggregation forgets; elimination remembers}
Keeping every component $P_a x$ is an exact change of representation when $\sum_aP_a=I$. Reconstruct by summing the components. In contrast, replacing $(p_0,p_1,p_2,p_3)$ by $(p_0+p_1,p_2+p_3)$ forgets how mass was distributed inside either pair. That information can matter for later observations.

Eliminating a hidden state is different again. Its coordinates disappear from the displayed state, but the exact retained equation acquires a history integral. Simply deleting the coupling produces the dashed curve above, not the original evolution.
\begin{center}\small
\begin{tabularx}{\linewidth}{@{}lYl@{}}\toprule
Question & Worked route & Plates\\\midrule
When is coarse evolution exact? & Aggregate columns, then test intertwining. & \plateRangeV{4}{6}\\
When is a coarse posterior exact? & Compare two hidden shapes with the same masses. & \plateRangeV{9}{13}\\
What changes under a fractional power? & Keep modes, change spectral rates. & \plateRangeV{15}{17}\\
What returns after elimination? & Compute the self-energy and memory kernel. & \plateRangeV{21}{25}\\\bottomrule
\end{tabularx}\end{center}
\takeaway{The chapter supplies exact identities and conditional bounds. Whether a particular coarse model satisfies their hypotheses is a separate calculation.}

\vpage{2}{A dictionary with orientations}{READ BEFORE MULTIPLYING / CONVENTIONS}{Chapter VI, Part I: Sections 2, 4, 6, 7, and 12. Column probabilities and backward observables.}
An apparent disagreement between two transfer formulas often comes from acting on different objects. Throughout this guide, probabilities are columns and observables are functions written as columns. Their pairing is $f^{\mathsf T}p$. The transpose distinguishes the forward action on distributions from the backward action on observations.
\begin{center}\small
\begin{tabularx}{\linewidth}{@{}lYY@{}}\toprule
Object & Exact meaning & What to check\\\midrule
$p$ & A nonnegative column with $\ones^{\mathsf T}p=1$. & It is not a vector of function values.\\
$Q$ & Continuous Markov generator. & Off-diagonal entries are nonnegative and columns sum to zero.\\
$K$ & One discrete stochastic step. & $\ones^{\mathsf T}K=\ones^{\mathsf T}$ and $K_{ij}\ge0$.\\
$K^{\mathsf T}f$ & Expected next observable given the current state. & The same dynamics, with the opposite action on the dual object.\\
$P_a$ & A coordinate packet projector. & All packets together reconstruct the full state.\\
$C$ & Sums fine probability masses into coarse blocks. & Usually rectangular and not invertible.\\
$E_\pi$ & Conditional expectation of observables. & A weighted average, self-adjoint for the $\pi$ inner product.\\
$\mathcal B_\lambda$ & $\operatorname{diag}(\lambda)p/(\lambda^{\mathsf T}p)$. & Nonnegative likelihood and positive total evidence.\\
$L$ & A nonnegative self-adjoint operator. & Its evolution is $e^{-tL}$, not $e^{tL}$.\\
$R(z)$ & $(zI-A)^{-1}$. & $z$ belongs to the resolvent set.\\
$\mathcal K(t)$ & The memory kernel after hidden-state elimination. & It is not the transition matrix $K$.\\\bottomrule
\end{tabularx}\end{center}
\head{One convention test}
For
\[
K_2=\begin{pmatrix}3/4&1/4\\1/4&3/4\end{pmatrix},
\qquad p=e_0,\qquad f=(2,1)^{\mathsf T},
\]
we obtain $K_2p=(3/4,1/4)^{\mathsf T}$ and
\[
f^{\mathsf T}K_2p=7/4=(K_2^{\mathsf T}f)_0.
\]
With departure weights $w_i=e^{-V(i)}$, the backward Feynman--Kac operator is $\operatorname{diag}(w)K_2^{\mathsf T}$. Its forward dual is $K_2\operatorname{diag}(w)$. Interchanging these factors changes when the weight is applied.
\takeaway{Always name the object being transported. The transpose and the position of the likelihood or potential are mathematical choices, not typographic details.}

\vpage{3}{A clock generated by one matrix}{CONTINUOUS TRANSFER / THE ENTIRE SOLUTION}{Chapter VI, Part I: Theorems 1.2 and 4.1. Finite matrix exponentials and Markov positivity.}
Take the generator
\[
G=\begin{pmatrix}-1&1\\1&-1\end{pmatrix}.
\]
The total-mass direction $(1,1)$ has eigenvalue zero. The contrast $(1,-1)$ has eigenvalue $-2$. Resolving those two directions gives the full transfer, not only its first-order approximation:
\[
T(t)=e^{tG}=\frac12\begin{pmatrix}1+e^{-2t}&1-e^{-2t}\\1-e^{-2t}&1+e^{-2t}\end{pmatrix}.
\]
\vfig{clock}
\captiontext{A point seed at state zero approaches the uniform distribution. The displayed times were selected so that every bar label is rational. This is continuous-time evolution, not a simulated trajectory.}
\head{Two checks explain the mechanism}
The constant mode is unchanged, while each elapsed interval multiplies the contrast by $e^{-2t}$. Consequently two intervals multiply their contrast factors:
\[
e^{-2s}e^{-2t}=e^{-2(s+t)},\qquad T(s)T(t)=T(s+t).
\]
For $t\ge0$, the entries are nonnegative and every column sums to one. Differentiation gives $T'(t)=GT(t)$. These three facts separately check the initial condition, the evolution equation, and preservation of probability.

At $t=\log 2$, the retained first coordinate is $5/8$. It would instead be $1/2$ under the isolated scalar equation $\dot x=-x$. The discrepancy is the mass that can leave state zero and return. On \pagerefV{22}, we remove the second coordinate correctly and recover exactly the same $5/8$ using a memory integral.
\takeaway{The semigroup lives on the complete state. One diagonal entry of a semigroup need not itself satisfy a scalar semigroup law.}

\vpage{4}{One evolution, two packet choices}{FOUR-STATE MODEL / COORDINATES AND MODES}{Chapter VI, Part I: Theorem 2.1, Corollary 2.2, and the packet path expansion.}
Our main finite model has generator
\[
Q=\begin{pmatrix}-3&1&2&0\\1&-3&0&2\\2&0&-3&1\\0&2&1&-3\end{pmatrix}.
\]
The graph shows rates between actual states. The four normalized sign vectors on the right form an orthonormal basis $u_j$. Their projectors $P_j=u_ju_j^{\mathsf T}$ sum to the identity and satisfy $P_iP_j=0$ for $i\ne j$.
\vfig{generator_modes}
\captiontext{State coordinates have off-diagonal transfers. Spectral packets reduce the generator and evolve independently. The two dashed state groups are the coarse partition used next.}
\head{The same evolution has a closed formula}
Writing $q=e^{-2t}$ gives
\[
e^{tQ}=P_0+qP_1+q^2P_2+q^3P_3.
\]
At $q=1/2$, the point seed $e_0$ becomes $(15,5,9,3)^{\mathsf T}/32$. The polynomial identity $T(q)T(r)=T(qr)$ checks composition exactly, without numerical matrix exponentiation.

Any other complete packet family still has exact block convolution. For the coordinate groups A and B,
\[
P_b e^{(s+t)Q}P_a=\sum_{c\in\{\mathrm A,\mathrm B\}}(P_b e^{sQ}P_c)(P_c e^{tQ}P_a).
\]
The intermediate packet cannot generally be omitted. Only a reducing family eliminates all cross terms.

For later discrete calculations define $K=I+Q/8$. This matrix is stochastic, but \emph{it is not} $e^{Q/8}$. It is the step kernel in the uniformization identity $e^{tQ}=e^{-8t}\sum_{n\ge0}(8t)^nK^n/n!$.
\takeaway{Changing packet coordinates does not change the evolution. Replacing continuous time by one discrete step does.}

\vpage{5}{When two states are truly equivalent}{EXACT LUMPABILITY / AGGREGATE COLUMNS}{Chapter VI, Part I: Theorem 5.1. The hypothesis is an intertwining identity.}
Group states $0,1$ into A and states $2,3$ into B. The aggregation map and the candidate coarse kernel are
\[
C=\begin{pmatrix}1&1&0&0\\0&0&1&1\end{pmatrix},\qquad
\bar K=\begin{pmatrix}3/4&1/4\\1/4&3/4\end{pmatrix}.
\]
A fine state in A has probability $1/4$ of entering B in one step, regardless of which fine state in A was occupied. The analogous statement holds for B.
\vfig{lumpability}
\captiontext{The fine columns are visibly different. Their sums over each target packet agree within each source packet. That equality is the reason a two-state description works.}
\head{A one-step calculation controls every discrete horizon}
The displayed columns give $CK=\bar KC$. Multiplication then yields
\[
CK^n=\bar K^nC\qquad(n\ge0).
\]
For example, $e_0$ and $e_1$ have the same coarse initial vector $(1,0)^{\mathsf T}$. Their fine evolutions differ, but their coarse vectors agree at every step. No choice of the within-A conditional distribution is needed for the prediction.

The continuous generator also closes:
\[
CQ=\bar Q C,\qquad
\bar Q=\begin{pmatrix}-2&2\\2&-2\end{pmatrix}.
\]
The Duhamel identity therefore gives $Ce^{tQ}=e^{t\bar Q}C$. This is a second consequence of the corresponding generator identity, not an inference that $K$ is a fixed-time sample of $Q$.

Aggregation is still not invertible. The theorem says the discarded details are irrelevant for \emph{these predictions}. It does not say they are irrelevant to every later likelihood. The observation counterexample on \pagerefV{10} uses this exact same partition.
\takeaway{A partition is not automatically a valid Markov model. The aggregate transition columns must be constant within each source block.}

\vpage{6}{A broken lump and its error budget}{DEFECTS / A GUARANTEE WITHOUT A FIT}{Chapter VI, Part I: Theorems 3.1, 3.3, and 5.2. Continuous and discrete defect identities.}
Modify only column zero of $K$: move $1/16$ of its mass from target state zero to target state two. The perturbed matrix $K_\varepsilon$ remains stochastic, but its aggregate column now differs from column one. Relative to the unchanged coarse kernel,
\[
\Delta=CK_\varepsilon-\bar KC
=\begin{pmatrix}-1/16&0&0&0\\1/16&0&0&0\end{pmatrix},
\qquad \|\Delta\|_1=1/8.
\]
\vfig{defect_budget}
\captiontext{The solid values are computed with rational matrix powers. The dashed certificate uses only stochasticity and the exact one-step defect. It need not be sharp.}
\head{Telescoping tells us where every error originates}
For $n\ge0$,
\[
CK_\varepsilon^n-\bar K^nC
=\sum_{j=0}^{n-1}\bar K^{n-1-j}\Delta K_\varepsilon^j.
\]
Each summand propagates one defect between an earlier fine evolution and a later coarse evolution. Column-stochastic matrices have induced $\ell^1$ norm one, so the norm is at most $n/8$. The empty sum at $n=0$ is zero.

The continuous counterpart is equally explicit. With fine generator $-1$, coarse generator $-2$, and $C=1$, the generator defect is $1$ and
\[
e^{-t}-e^{-2t}=\int_0^t e^{-2(t-r)}e^{-r}\,dr.
\]
Its resolvent version is $1/(z+1)-1/(z+2)=1/((z+1)(z+2))$ for $z>0$.
\takeaway{A small defect does not disappear. The identity specifies how it accumulates, and the norm assumptions turn that identity into a finite-horizon certificate.}

\vpage{7}{Averaging is not summing}{OBSERVABLES / WEIGHTED ORTHOGONALITY}{Chapter VI, Part I: Theorem 6.1. Strictly positive weights are part of the setting.}
For this observable example, choose the fixed probability weights $\pi=(1,2,3,4)/10$. They are not being identified with the stationary law of the earlier symmetric chain. On the same two-block partition, conditional expectation is
\[
E=\begin{pmatrix}
1/3&2/3&0&0\\1/3&2/3&0&0\\0&0&3/7&4/7\\0&0&3/7&4/7
\end{pmatrix}.
\]
It replaces function values by their \emph{weighted block average}. In particular, $f=(0,3,2,5)^{\mathsf T}$ becomes $Ef=(2,2,26/7,26/7)^{\mathsf T}$.
\vfig{weighted_projection}
\captiontext{The bars are function values, not state probabilities. Their total need not be one. The block masses $3/10$ and $7/10$ determine the denominators in the averaging formula.}
\head{Which inner product makes this an orthogonal projection?}
Put $W=\operatorname{diag}(\pi)$. Direct multiplication checks
\[
E^2=E,\qquad WE=E^{\mathsf T}W,\qquad E\ones=\ones.
\]
The middle equation is weighted self-adjointness. Ordinary symmetry fails: $E_{01}=2/3$ whereas $E_{10}=1/3$. Testing $E=E^{\mathsf T}$ would therefore reject a correct conditional-expectation projector.

The range consists exactly of functions constant on each block. The error $f-Ef$ has weighted mean zero inside every block, so it is weighted-orthogonal to every block-constant function. This is the finite projection interpretation of conditional expectation.

In contrast, $C$ acts on probability columns by summing masses. The shapes of $C$ and $E$ differ, and their domains have different interpretations. The notation ``coarsening'' alone does not specify which operation is intended.
\takeaway{An exact weighted projector can be oblique in the unweighted Euclidean picture. Keep the measure and the inner product in the model.}

\vpage{8}{Where the variance goes}{MARTINGALE PACKETS / THREE EXACT PIECES}{Chapter VI, Part I: Theorem 6.2. Nested weighted projections and their differences.}
Keep the weights and observable from \pagerefV{7}. Let $E_0f$ be the constant global mean, let $E_1=E$ average within A and B, and let $E_2=I$ retain every value. Define
\[
D_0=E_0,\qquad D_1=E_1-E_0,\qquad D_2=I-E_1.
\]
Then $D_0f$ is the mean, $D_1f$ describes the contrast between block means, and $D_2f$ is the fluctuation left inside each block.
\vfig{variance}
\captiontext{The strip lengths are proportional to the three exact weighted energies. The between-block component is small but not zero. The vectors below specify the corresponding directions, not only their sizes.}
\head{The energy identity follows from actual orthogonality}
The tower identities give $E_0E_1=E_1E_0=E_0$. Together with weighted self-adjointness, these imply
\[
D_iD_j=0\ (i\ne j),\qquad \sum_{j=0}^2 D_j=I.
\]
For the chosen $f$, the total weighted energy is $\|f\|_\pi^2=13$ and the mean is $16/5$. Consequently
\[
\operatorname{Var}_\pi(f)
=13-\left(\frac{16}{5}\right)^2
=\frac{69}{25}
=\frac{108}{175}+\frac{15}{7}.
\]
The first term is the energy of the between-block contrast, and the second is the weighted within-block variance. This is a decomposition of one observable in one fixed weighted space.

Changing the probability weights between levels would change the inner products and the averaging operators. The displayed orthogonal decomposition would no longer follow merely from nesting the sets. Likewise, choosing unrelated projections can preserve the telescope $\sum D_j=I$ while losing orthogonality and the energy sum.
\takeaway{The differences between nested conditional expectations are not just bookkeeping. They are orthogonal, separately measurable parts of the original signal.}

\vpage{9}{Evidence changes shape and mass}{BAYESIAN TRANSFER / AN EXACT POSTERIOR}{Chapter VI, Part I: Theorems 7.1--7.3. Normalization, packet reconstruction, and positive evidence.}
Start with $p=(1,1,1,1)^{\mathsf T}/4$ and likelihood weights $\lambda=(4,1,1,1)^{\mathsf T}$. The unnormalized vector is $(1,1/4,1/4,1/4)^{\mathsf T}$, its evidence is $Z=7/4$, and
\[
\mathcal B_\lambda(p)=\frac{(4,1,1,1)^{\mathsf T}}7.
\]
Likelihood weights may be rescaled, so an evidence larger than one is not a probability claim. Multiplying all likelihoods by one common positive factor leaves this posterior unchanged.
\vfig{evidence}
\captiontext{The A packet becomes heavier and its internal shape changes. The B packet retains an equal internal split but loses total mass. These are different effects of the same likelihood.}
\head{Reconstruct from packet mass times packet shape}
The unnormalized packet evidences are $Z_{\mathrm A}=5/4$ and $Z_{\mathrm B}=1/2$. The resulting masses and conditional shapes are
\[
m_{\mathrm A}=5/7,\quad q_{\mathrm A}=(4/5,1/5,0,0)^{\mathsf T},
\qquad m_{\mathrm B}=2/7,\quad q_{\mathrm B}=(0,0,1/2,1/2)^{\mathsf T}.
\]
Thus $m_{\mathrm A}q_{\mathrm A}+m_{\mathrm B}q_{\mathrm B}$ is exactly the displayed posterior. Keeping only the masses would not reconstruct the fine state.

There is no need to divide by a zero packet evidence. For likelihood $(0,0,1,1)$, packet A has mass zero and contributes the zero vector. Only the occupied B packet has a conditional probability vector. If the \emph{total} evidence is zero, however, the Bayesian update is not defined and the implementation rejects the input.
\takeaway{Bayesian reconstruction uses two kinds of data: how much posterior mass a packet carries and how that mass is distributed inside the packet.}

\vpage{10}{Two priors that fool the coarse model}{SUFFICIENCY / A COMPUTABLE FAILURE WITNESS}{Chapter VI, Part I: Theorem 8.2 and the one-packet exception.}
Keep the same likelihood $\lambda=(4,1,1,1)^{\mathsf T}$ and partition A$=\{0,1\}$, B$=\{2,3\}$. Compare
\[
p^{(a)}=(1/2,0,1/2,0)^{\mathsf T},\qquad
p^{(b)}=(0,1/2,1/2,0)^{\mathsf T}.
\]
Both have coarse prior $(1/2,1/2)^{\mathsf T}$. A procedure given only that two-entry vector cannot distinguish them.
\vfig{insufficiency}
\captiontext{After the identical fine likelihood is applied, the coarse posteriors are $(4/5,1/5)$ and $(1/2,1/2)$. A single deterministic update of the common coarse prior cannot produce both.}
\head{Why this witnesses the necessity of blockwise constancy}
The A evidence for $p^{(a)}$ is $2$, whereas it is $1/2$ for $p^{(b)}$. Their B evidence is $1/2$ in both cases. The posterior odds therefore distinguish the hidden distribution inside A.

For strictly positive fine and coarse likelihoods, and at least two coarse blocks, the manuscript proves the all-prior criterion
\[
C\mathcal B_\lambda(p)=\mathcal B_{\bar\lambda}(Cp)\quad\hbox{for every }p
\quad\Longleftrightarrow\quad
\lambda_i=c\bar\lambda_a\quad(i\in I_a)
\]
for one global $c>0$. The global scale is invisible after normalization. A different arbitrary scale for each packet is not invisible: it changes posterior odds.

If all states are merged into a single packet, every normalized posterior coarsens to $(1)$. That degenerate commutation holds for every likelihood and cannot identify any within-packet information. It is explicitly excluded from the nontrivial criterion.
\takeaway{A prior-dependent effective likelihood can always summarize the update. A fixed likelihood that works for every prior is a much stronger requirement.}

\vpage{11}{The exact information-loss vector}{BAYESIAN DEFECT / MORE THAN AN ERROR BAR}{Chapter VI, Part I: Theorems 8.1 and 8.4, and Corollary 8.6.}
For a prior $p$, let $P=Cp$ and let $r_a(p)$ be the conditional average of the fine likelihood inside packet $a$. For the uniform prior and $\lambda=(4,1,1,1)$,
\[
P=(1/2,1/2)^{\mathsf T},\qquad r=(5/2,1)^{\mathsf T}.
\]
The exact coarse posterior is $\mathcal B_r(P)=(5/7,2/7)^{\mathsf T}$. Suppose instead that the declared coarse likelihood is $\bar\lambda=(1,1)$.
\vfig{loss_vector}
\captiontext{The posterior error lies along the mass-preserving direction of the two-state simplex. Its entries are opposite because both endpoints have total mass one.}
\head{Compute the factorization, including the null direction}
Here $\bar b=\mathcal B_{\bar\lambda}(P)=P$, $\delta=r-\bar\lambda=(3/2,0)^{\mathsf T}$, and $Z_r=7/4$. Hence
\[
\frac1{Z_r}\left(\diag(P)-\bar bP^{\mathsf T}\right)\delta
=\frac4{7}\begin{pmatrix}1/4&-1/4\\-1/4&1/4\end{pmatrix}
\begin{pmatrix}3/2\\0\end{pmatrix}
=\begin{pmatrix}3/14\\-3/14\end{pmatrix}.
\]
This is the actual fine-to-coarse discrepancy, not a fitted approximation. The matrix kills $\bar\lambda$, recording the common-scale null direction. Its $\ell^1$ error is $3/7$.

For a family of likelihoods, states may be grouped by equality of their \emph{whole likelihood profile}. The family $(3,3,1,1)$ and $(1,2,1,2)$ separates all four states even though either likelihood alone permits a coarser grouping. The common refinement is therefore the discrete partition. The theorem's one-packet exception still remains.
\takeaway{The defect factors through exactly the conditional evidence information missing from the declared coarse model. A new observation family can require a finer partition.}

\vpage{12}{A nonlinear semigroup of odds}{BAYESIAN FLOW / COMPOSITION WITHOUT LINEARITY}{Chapter VI, Part I: Theorems 7.1 and 9.1. Likelihood composition and the replicator field.}
Let the score vector be $g=(1,0)$ and start at $(1/2,1/2)$. The likelihood at time $t$ is $(e^t,1)$, so the first posterior coordinate is
\[
p_0(t)=\frac{e^t}{1+e^t},\qquad p_1(t)=\frac1{1+e^t}.
\]
The odds $p_0/p_1$ grow exponentially, while the probabilities remain inside the simplex. At $t=\log 2$ the posterior is exactly $(2/3,1/3)$.
\vfig{odds}
\captiontext{The solid and dashed curves sum to one. The curvature is caused by the shared evidence denominator, not by a nonlinear likelihood multiplication law.}
\head{Two properties that must not be conflated}
If $a>0$ multiplies the odds, the first-coordinate map is
\[
\phi_a(x)=\frac{ax}{1+(a-1)x}.
\]
Direct algebra gives $\phi_a(\phi_b(x))=\phi_{ab}(x)$. Thus $a=e^t$ gives a nonlinear semigroup. Nevertheless, $\phi_3(1/2)=3/4$, while the mean of $\phi_3(0)$ and $\phi_3(1)$ is $1/2$. Superposition fails.

Differentiation of the normalized flow yields the centered score equation
\[
\dot p_i=\bigl(g_i-\langle g,p\rangle\bigr)p_i.
\]
For the displayed model, $\dot p_0=p_0(1-p_0)$. The negative average-score term is what keeps the total derivative zero. Omitting the derivative of the normalizer would not preserve mass.

Likelihood composition also preserves the evidence identity: the first evidence times the second conditional evidence equals the combined evidence. Normalization cancels in the posterior, but that factor still matters when comparing data models.
\takeaway{A semigroup is a composition law in time. It does not imply that the maps are linear or that they preserve mixtures.}

\vpage{13}{Prediction and observation do not commute}{HIDDEN MARKOV FILTERING / TIME ORDER}{Chapter VI, Part I: Proposition 9.3. Linear prediction followed by a positive-evidence update.}
Use the two-state kernel $K_2$ from \pagerefV{2}, the initial distribution $e_0$, and likelihood $(1,1/2)$. Predicting first gives $(3/4,1/4)$. Its likelihood-weighted vector is $(3/4,1/8)$, with evidence $7/8$, so the posterior is $(6/7,1/7)$.

Updating first leaves $e_0$ unchanged. Predicting afterward gives $(3/4,1/4)$ instead. The order of operations expresses which time the observation concerns.
\vfig{filter_order}
\captiontext{Both routes produce valid probability vectors. Their difference is therefore not a normalization bug. It is a different conditioning problem.}
\head{Keep an unnormalized path alongside the filter}
For a sequence of observations, let $D_r=\diag(\lambda^{(r)})$. The forward recursion is
\[
\widehat p_r=Kp_{r-1},\qquad p_r=D_r\widehat p_r/Z_r.
\]
Its unnormalized companion is
\[
v_r=D_rK\cdots D_1Kp_0,
\qquad p_r=\frac{v_r}{\ones^{\mathsf T}v_r},
\qquad \ones^{\mathsf T}v_r=\prod_{j=1}^r Z_j.
\]
The product of evidences records the normalization lost at each individual step. The four-state computation on \pagerefV{26} checks both paths independently, including their coarse versions.

Here $D_rK$ means an observation after prediction. In the Feynman--Kac example on \pagerefV{20}, departure-state weighting has forward operator $K D_w$. Both conventions are legitimate once their meanings are fixed, but they cannot be substituted for one another.
\takeaway{Write the chronological sequence before multiplying matrices. A product can be normalized correctly and still answer the wrong time-ordered question.}

\vpage{14}{Conjugacy tracks exponents}{BETA MODEL / POSTERIOR AND EVIDENCE}{Chapter VI, Part I: Example 9.2. A finite-parameter conjugate family, not a universal posterior model.}
For a small exact example, take the $\mathrm{Beta}(2,3)$ density
\[
\rho(\theta)=12\theta(1-\theta)^2,\qquad0\le\theta\le1.
\]
Observe a specified ordered sequence with two successes and one failure. Its likelihood is $\theta^2(1-\theta)$. Multiplication adds exponents, so the posterior belongs to $\mathrm{Beta}(4,4)$ with density $140\theta^3(1-\theta)^3$.
\vfig{beta}
\captiontext{The dashed prior is asymmetric. The solid posterior is symmetric with mean $1/2$. Polynomial integration verifies the normalizers exactly; the curves are sampled drawings of those polynomials.}
\head{The observation convention changes the evidence}
The integral of the unnormalized posterior is
\[
\int_0^1 12\theta^3(1-\theta)^3\,d\theta=\frac3{35}.
\]
Dividing by $3/35$ gives the displayed normalized density. If the observation is instead the \emph{count} of two successes in three trials, there are $\binom32=3$ compatible ordered sequences. The count likelihood is therefore $3\theta^2(1-\theta)$ and its evidence is $9/35$.

The extra factor does not depend on $\theta$, so it cancels from the posterior. The two experiments thus have the same posterior distribution but different data probabilities. Evidence is not determined by the final normalized density alone.

The exact check also integrates both densities to one and integrates $\theta$ against the posterior to obtain $1/2$. This demonstrates one conjugate family in elementary terms. It does not show that arbitrary likelihoods preserve that family or that all posterior information can be compressed into two exponents.
\takeaway{Conjugacy gives a closed parameter update when a specified likelihood respects a specified family. Preserve the evidence as well as the normalized posterior.}

\vpage{15}{The modes stay; their clocks change}{FRACTIONAL DIFFUSION / NEW COUPLINGS}{Chapter VI, Part I: Theorem 10.1. A positive self-adjoint finite model and its spectral powers.}
Return to the four-state generator on \pagerefV{4} and put $L=-Q$. It is real symmetric and nonnegative, with orthogonal projectors $P_0,P_1,P_2,P_3$ and eigenvalues $0,2,4,6$. Its square root is
\[
L^{1/2}=\sqrt2P_1+2P_2+\sqrt6P_3,
\qquad (L^{1/2})^2=L.
\]
Both heat evolutions retain these modes. Ordinary diffusion multiplies mode $j$ by $e^{-2jt}$; square-root diffusion multiplies it by $e^{-\sqrt{2j}\,t}$. The constant mode is unchanged in both.
\vfig{fractional}
\captiontext{The highest mode decays at rates $6$ and $\sqrt6$. The graph beside the curves shows the additional couplings in the fractional operator. Curves display the explicit formulas, not an empirical fit.}
\head{Fractional does not mean entrywise}
Although $L_{03}=0$, its spectral square root satisfies
\[
(L^{1/2})_{03}=\frac{\sqrt6-\sqrt2-2}{4}<0.
\]
Thus $-L^{1/2}$ permits a direct transition between states $0$ and $3$, which were not neighbors for $Q$. Taking the square root of each matrix entry would neither produce this value nor square back to $L$.

The exact block aggregation remains compatible in this example:
\[
C L^{1/2}=\mat{1&-1\\-1&1}C.
\]
The coarse matrix is the positive square root of $-Q_c$. This is a concrete naturality identity, checked on the full matrices. It is not a claim that an arbitrary partition commutes with every fractional construction.

The computation uses a nonnegative self-adjoint operator and a specified spectral branch. General nonnormal matrices, unbounded domains, and positivity of other fractional evolutions need their own hypotheses.
\takeaway{Fractional powers alter modal decay rates. Reassembling those modes can introduce long-range spatial interactions even when the original generator was sparse.}

\vpage{16}{A gap changes the error scale}{FRACTIONAL DEFECTS / NEAR ZERO VERSUS AWAY FROM ZERO}{Chapter VI, Part I: Theorem 10.3 and Corollary 10.4. Separate generator error from fractional error.}
The function $a\mapsto\sqrt a$ changes most rapidly near zero. A scalar model already shows why a spectral gap appears in fractional-defect estimates:
\[
|\sqrt\varepsilon-\sqrt0|=\sqrt\varepsilon,
\qquad
\sqrt{1+\varepsilon}-1
=\frac{\varepsilon}{\sqrt{1+\varepsilon}+1}\le\frac\varepsilon2.
\]
The same input difference has a square-root scale at zero and a linear bound above a fixed positive lower bound.
\vfig{gap}
\captiontext{The upper curve compares $0$ with $\varepsilon$. The lower curve compares $1$ with $1+\varepsilon$. The dotted line $\varepsilon/2$ bounds the latter. All plotted values come from the displayed scalar formulas.}
\head{Read an operator bound with its hypotheses}
For positive self-adjoint finite matrices $A,B\ge mI$ and $0<\alpha<1$, the manuscript's gapped estimate has the form
\[
\norm{CA^\alpha-B^\alpha C}
\le \alpha m^{\alpha-1}\norm{CA-BC}.
\]
It concerns an intertwining defect, not just two eigenvalue lists. With $C=1$, $A=4$, $B=9$, $m=4$, and $\alpha=1/2$, the generator defect has size $5$, the fractional defect has size $1$, and the bound is $5/4$.

The common zero mode of a Markov generator cannot simply be included in a positive-gap estimate. One must either work on an appropriate gapped subspace, handle the zero mode separately, or use the gapless estimate. In particular, a good estimate on each finite level need not have a constant uniform in the refinement level.

On \pagerefV{27}, the sequences $1/n^2$ and $1/n$ turn this change in error scale into a concrete distinction between summable and nonsummable defects.
\takeaway{A fixed gap is quantitative information. It prevents the derivative of the fractional function from becoming singular along the comparison.}

\vpage{17}{A random diffusion clock}{SUBORDINATION / A DISTRIBUTION OF ORDINARY TIMES}{Chapter VI, Part I: Theorem 10.2. The Poisson subordination formula for the finite positive model.}
Square-root evolution can be written as an average of ordinary diffusion over a positive clock variable $s$:
\[
e^{-t\sqrt L}=\int_0^\infty \eta_t(s)e^{-sL}\,ds,
\qquad
\eta_t(s)=\frac{t}{2\sqrt\pi}s^{-3/2}e^{-t^2/(4s)},\quad t>0.
\]
Here $t$ is the external time; $s$ is the ordinary diffusion time being mixed. They are not interchangeable variables.
\vfig{random_clock}
\captiontext{The density $\eta_1$ peaks at $s=1/6$ but has a long right tail. The drawing ends at $s=2$; it does not claim that the missing tail has zero mass.}
\head{Check the normalizer before interpreting a density}
Under $s=t^2/(4u^2)$, reversing the integration limits gives
\[
\eta_t(s)\,|ds|=\frac2{\sqrt\pi}e^{-u^2}\,du.
\]
Its integral over $u>0$ is one, so $\eta_t$ is a probability density. On each eigenmode of $L$, the operator formula reduces to the scalar Laplace identity
\[
\int_0^\infty\eta_t(s)e^{-\lambda s}\,ds=e^{-t\sqrt\lambda}.
\]
The manuscript supplies that scalar identity and its spectral assembly. The appendix separately checks the change of variables and normalization; sampled curves are not used to certify an improper integral.

A normalized density need not have a finite mean. For large $s$, the exponential factor tends to one and $\eta_t(s)$ is proportional to $s^{-3/2}$. Consequently $s\eta_t(s)$ has a nonintegrable $s^{-1/2}$ tail. The clock has infinite mean even though its total mass is one.
\takeaway{Subordination is a weighted mixture of diffusion times. Replacing the clock distribution by a single average time is not justified, and here that average is infinite.}

\vpage{18}{A Green kernel is an inverse with units}{DIRICHLET GRID / THE MISSING FACTOR OF h}{Chapter VI, Part I: Theorem 11.1. Exact nodal Green reconstruction for this one-dimensional stencil.}
Take three interior nodes $x_i=i/4$ on $[0,1]$, with zero boundary values and $h=1/4$. The discrete negative second derivative is
\[
L_h=\frac1{h^2}\mat{2&-1&0\\-1&2&-1\\0&-1&2}.
\]
For the continuous Dirichlet problem, the Green function is $G(x,y)=\min(x,y)-xy$. At these nodes the exact comparison is
\[
L_h^{-1}=h\,[G(x_i,x_j)]_{i,j=1}^3
=\frac1{64}\mat{3&2&1\\2&4&2\\1&2&3}.
\]
\vfig{green}
\captiontext{The tent is the central column of the continuous Green kernel. It has nodal values $(1,2,1)/8$; its discrete inverse column is $(1,2,1)/32$, smaller by the factor $h$.}
\head{A forcing vector is not a quadrature-weighted measure}
For sampled forcing $f_j$, the finite solution is
\[
u_i=\sum_j(L_h^{-1})_{ij}f_j
=h\sum_jG(x_i,x_j)f_j.
\]
The factor $h$ converts samples into the quadrature weights appropriate to this formula. Dropping it multiplies the computed solution by four on the displayed grid.

A discrete unit coordinate vector and a discretized unit point mass use different normalizations. For $f=e_2$, the result is the second inverse column. For $f=e_2/h$, it is the unscaled tent sampled at the nodes. Either convention can be useful, but it must be specified before comparing answers.

Exact checks multiply $L_h$ by the proposed inverse at $1,3,7,$ and $15$ interior nodes. The general identity is the manuscript theorem. These finite checks do not extend it to arbitrary meshes, dimensions, or boundary conditions.
\takeaway{An inverse kernel carries a measure convention. A correct-looking tent can give a wrong solution if the grid weight is omitted.}

\vpage{19}{Convergence with a nonzero exact error}{FINITE DIFFERENCES / A QUARTIC TEST PROBLEM}{Chapter VI, Part I: Theorems 11.1--11.2. Exact nodal inversion and the scale of low-mode errors.}
Use $u(x)=x-x^4$, which vanishes at both endpoints and solves $-u''=12x^2$. Its centered second difference is not exact:
\[
L_hu|_{\rm nodes}=12x_i^2+2h^2.
\]
Since the stencil differentiates the quadratic $x(1-x)$ exactly, the solution $u_h$ to the sampled forcing $12x_i^2$ satisfies
\[
\boxed{u_h(x_i)-u(x_i)=-h^2x_i(1-x_i).}
\]
This gives both the sign and size of the error, rather than merely a log--log slope.
\vfig{mesh}
\captiontext{The three curves show magnitudes of the exact nodal errors for $h=1/4,1/8,1/16$. Lines between grid points are visual guides, not an additional continuous numerical solution. The error is negative on every interior node.}
\head{One refinement divides this error by four}
All three grids include $x=1/2$, where the maximum absolute nodal error is $h^2/4$. It equals $1/64$, $1/256$, and $1/1024$, respectively. Exact rational matrix inversion verifies these values independently of the closed error formula.

The manuscript also compares the discrete eigenvalues
\[
\lambda_{k,h}=4h^{-2}\sin^2(k\pi h/2)
\]
with $(k\pi)^2$, with error at most $(k\pi)^4h^2/12$. For a fixed mode $k$, the bound tends to zero quadratically. It does not claim a uniform small error for modes whose index grows like $1/h$.

The nodal Green identity is exact, while the discretization of the differential equation still has a truncation error. These statements concern different comparisons and do not contradict one another.
\takeaway{An exact inverse of the discrete operator does not make that operator equal to the differential operator. Keep algebraic inversion and continuum approximation separate.}

\vpage{20}{Eight paths, one weighted expectation}{FEYNMAN--KAC / DEPARTURE-STATE WEIGHTS}{Chapter VI, Part I: Theorem 12.1. A finite path sum with its time convention displayed.}
Let $K_2=\frac14\mat{3&1\\1&3}$, start at state $0$, use weights $w=(1,1/2)$, and terminal observable $f=(2,1)$. For three transitions, there are eight possible state sequences after the fixed initial state. We evaluate
\[
\EE_0\left[\prod_{r=0}^{2}w(X_r)\,f(X_3)\right].
\]
Each path contributes its transition probability times its three departure weights times its terminal value.
\vfig{paths}
\captiontext{Every possible path is included. The probability column sums to one. The contribution column sums to $311/256$, agreeing exactly with the transfer calculation.}
\head{Check which side of the transition gets weighted}
The backward operator on observables is $D_w K_2^{\mathsf T}$. Therefore the expectation is the zeroth coordinate of
\[
(D_wK_2^{\mathsf T})^3f.
\]
Its forward dual is $K_2D_w$: weight the mass at the departing state, then move it. Weighting arrival states instead gives $D_wK_2$, a different expression. The exact tests deliberately check that the two conventions do not agree in this example.

The answer $311/256$ exceeds one, but it is not a probability. The terminal observable can equal two. With $0\le w_i\le1$, the weighted expectation is bounded by the maximum terminal value, not necessarily by one.

This calculation certifies a finite three-step identity for a stated model. The continuous-time theorem uses the corresponding infinitesimal generator and potential, with the direction of action kept consistent.
\takeaway{A Feynman--Kac weight changes the value assigned to each path. Its location relative to a transition is part of the model, not a typographical choice.}

\vpage{21}{Elimination makes a frequency-dependent term}{SCHUR--FESHBACH / RETAINED RESOLVENT}{Chapter VI, Part I: Theorem 13.1. A two-state Schur complement with the inverse checked directly.}
The two-state evolution from \pagerefV{3} has generator
\[
A=\mat{-1&1\\1&-1}.
\]
Retain the first coordinate and eliminate the second. For $z$ away from the relevant poles, the retained entry of $(zI-A)^{-1}$ is
\[
R_x(z)=\left[z+1-\frac1{z+1}\right]^{-1}
=\frac{z+1}{z(z+2)}.
\]
The term $\Sigma(z)=1/(z+1)$ is the hidden coordinate's self-energy.
\vfig{resolvent}
\captiontext{The full inverse gives a rational response with two poles, at $0$ and $-2$. Its partial fractions are $R_x(z)=\frac12z^{-1}+\frac12(z+2)^{-1}$. The hidden intermediate inverse has its own excluded point $z=-1$.}
\head{An eliminated variable has not disappeared}
If the coupling to the second coordinate were simply deleted, the retained resolvent would be $(z+1)^{-1}$. That has a different pole and a different long-time limit. The correction is frequency dependent, so it cannot be represented by one constant scalar generator on the retained coordinate.

The Schur expression first requires $z+1\ne0$, as well as invertibility of the full matrix. Its simplified rational form has a removable intermediate singularity at $z=-1$, where the retained entry of the full inverse is well-defined and equals zero. Simplification does not retroactively make the eliminated block invertible there.

For $z>0$, the Laplace transform of $(1+e^{-2t})/2$ is exactly $R_x(z)$. The next page derives the time-domain equation producing that response.
\takeaway{The Schur complement is exact elimination, not independent dynamics on a smaller state space. Hidden evolution returns as a $z$-dependent self-energy.}

\vpage{22}{Hidden state returns as memory}{EXACT RETAINED DYNAMICS / NOT A SCALAR SEMIGROUP}{Chapter VI, Part I: Theorem 13.2. The retained-state memory identity, including hidden initial data.}
Write the same system as
\[
x'=-x+y,\qquad y'=x-y.
\]
Solving the second equation first gives
\[
y(t)=e^{-t}y(0)+\int_0^t e^{-(t-r)}x(r)\,dr.
\]
Substitution is exact:
\[
x'(t)=-x(t)+e^{-t}y(0)+\int_0^t e^{-(t-r)}x(r)\,dr.
\]
The hidden initial state produces a separate forcing term. It must be retained even when the memory kernel is known.
\vfig{memory}
\captiontext{For $(x(0),y(0))=(1,0)$, the exact retained coordinate is $(1+e^{-2t})/2$. The dashed curve $e^{-t}$ solves the equation obtained by deleting the hidden coordinate rather than eliminating it.}
\head{Projection can destroy the semigroup law}
Let $R(t)=(1+e^{-2t})/2$, the response to this fixed initial preparation. At $t=\log2$,
\[
R(t)=\frac58,\qquad R(2t)=\frac{17}{32},
\qquad R(2t)-R(t)^2=\frac9{64}.
\]
The full two-state propagator is a semigroup, but this retained scalar response is not. Restarting it as a scalar model silently resets the hidden coordinate, which generally carries information from the first interval.

The appendix verifies the full differential equation, initial condition, integral memory equation, and Laplace resolvent in this model. Its kernel is $e^{-t}$, whose transform is the self-energy on \pagerefV{21}.
\takeaway{Exact reduction retains a record of the eliminated state through convolution and initial forcing. A memoryless approximation must justify the information it removes.}

\vpage{23}{A repeated pole remembers a Jordan chain}{NONDIAGONAL HIDDEN DYNAMICS / POLYNOMIAL EXPONENTIALS}{Chapter VI, Part I: Theorem 13.2 and Section 15. Functional calculus must retain nilpotent data.}
Now use a two-dimensional hidden block
\[
J=-2I+N,\qquad N=\mat{0&1\\0&0},\qquad N^2=0.
\]
Choose the entrance column $V=(0,1)^{\mathsf T}$ and exit row $U=(1,0)$. The hidden exponential is not $e^{-2t}I$:
\[
e^{tJ}=e^{-2t}(I+tN),
\qquad K_{\rm mem}(t)=Ue^{tJ}V=t e^{-2t}.
\]
The nilpotent link carries an input from the second hidden coordinate to the first.
\vfig{jordan}
\captiontext{The memory kernel starts at zero, becomes positive, and peaks at $t=1/2$. A rule that retained the eigenvalue $-2$ but discarded the Jordan chain would predict the zero kernel for these entrance and exit maps.}
\head{The transfer function sees pole order}
The exact transform is
\[
U(zI-J)^{-1}V=\frac1{(z+2)^2}
=\int_0^\infty t e^{-(z+2)t}\,dt,
\qquad \operatorname{Re}z>-2.
\]
The repeated pole and the polynomial factor $t$ are the same information in two domains. For a longer nilpotent chain, higher powers of $t$ occur with factorial denominators.

A zero value $K_{\rm mem}(0)$ does not make the entire kernel zero. Nor do stable eigenvalues alone specify every transient contribution: entrance maps, exit maps, and nilpotent structure matter. Here the full kernel is nevertheless elementary, so its peak and transform can be checked without numerical fitting.

This page is an explicit finite realization. It does not show that an arbitrary continuum memory kernel admits a finite-dimensional representation of this form.
\takeaway{Repeated eigenvalues may carry directional information through Jordan chains. Keep the finite jet of the functional calculus, not just its values on the spectrum.}

\vpage{24}{A zero initial hazard can still lead to certain exit}{PHASE-TYPE SURVIVAL / TWO SUCCESSIVE STAGES}{Chapter VI, Part I: Theorem 14.1. Survival, density, hazard, and moments have distinct meanings.}
Let a transient process move through two successive rate-one stages, then leave the system. With column probabilities, the transient generator and starting vector are
\[
T=\mat{-1&0\\1&-1},\qquad p(0)=e_0.
\]
Thus $p(t)=e^{-t}(1,t)^{\mathsf T}$. Its total remaining mass, exit density, and hazard are
\[
S(t)=(1+t)e^{-t},\qquad f(t)=t e^{-t},
\qquad h(t)=\frac{t}{1+t}.
\]
The initial hazard is zero because an exit needs two stages, not one.
\vfig{survival}
\captiontext{Survival falls from one to zero, while absorbed mass $1-S$ rises to one. The dotted hazard is an instantaneous rate, not a probability, even though it shares the displayed numerical scale. Time is measured in inverse stage-rate units.}
\head{Missing transient mass is absorbed mass}
The column sums of $T$ are $0$ and $-1$. The exit-rate row is therefore $(0,1)$, and $f(t)=(0,1)p(t)$. Adding an absorbing state restores a three-state probability vector with total mass one.

Exact integration gives
\[
\int_0^\infty f(t)\,dt=1,\qquad
\EE[\tau]=2,\qquad \EE[\tau^2]=6,
\qquad \Var(\tau)=2.
\]
Equivalently, the mean is $\ones^{\mathsf T}(-T)^{-1}e_0$. A single exponential waiting time with mean two would have constant hazard $1/2$ and variance four, so matching the mean alone would not reproduce this distribution.

The hazard formula is used where $S(t)>0$. The model has that property for every finite $t\ge0$. No division by an already vanished survival mass is required.
\takeaway{Survival tracks mass, density tracks exits per unit time, and hazard conditions on survival. Equal means or equal initial rates do not identify a waiting-time law.}

\vpage{25}{Naturality also transports derivatives}{HOLOMORPHIC CALCULUS / A RECTANGULAR COMPARISON}{Chapter VI, Part I: Theorem 15.1. Exact intertwining passes to the specified functional calculus.}
Consider a fine operator, a coarse operator, and a rectangular map
\[
A_f=\mat{2&1&0\\0&2&0\\0&0&9},\qquad
A_c=\mat{2&1\\0&2},\qquad
C_r=\mat{1&0&0\\0&1&0}.
\]
Direct multiplication gives $C_rA_f=A_cC_r$. For $F(z)=1/(1+z)$, which is defined at both spectra, compute the matrices by inversion:
\[
F(A_f)=\mat{1/3&-1/9&0\\0&1/3&0\\0&0&1/10}.
\]
\vfig{naturality}
\captiontext{The discarded eigenvalue $9$ contributes $1/10$ only on the fine side. On the retained Jordan block, the off-diagonal entry is $F'(2)=-1/9$, not zero. The rectangular comparison preserves that entry.}
\head{Why spectral values are not enough}
On the retained block $A_c=2I+N$, with $N^2=0$, one has
\[
F(A_c)=F(2)I+F'(2)N=\tfrac13I-\tfrac19N.
\]
Hence $C_rF(A_f)=F(A_c)C_r$. Replacing the block by $F(2)I$ would lose the nonzero derivative term and no longer give the actual inverse of $I+A_c$.

For polynomials, naturality follows by transporting every operator power. The manuscript's finite holomorphic statement uses finite Taylor data at each primary block and the corresponding nilpotent orders. This example verifies a rational instance with an actual inverse rather than asserting that arbitrary scalar substitutions are meaningful.

The same mechanism explains why the Jordan memory kernel on \pagerefV{23} contains $t e^{-2t}$. Exponentiation also has nonzero derivatives at repeated spectral values.
\takeaway{A functional calculus acts on operators, including their nilpotent parts. Exact intertwining transports that full action, not merely a list of scalar eigenvalue images.}

\vpage{26}{One complete coarse filter}{END-TO-END MODEL / TWO OBSERVATIONS}{Chapter VI, Part I: Proposition 9.3 with Theorems 5.1 and 8.2. Markov closure and likelihood sufficiency are separate.}
Use the four-state kernel $K$ and aggregation $C$ from \pagerefV{5}, starting at $p_0=e_0$. Observe two likelihood vectors constant inside both packets:
\[
\lambda^{(1)}=(3/4,3/4,1/4,1/4),\qquad
\lambda^{(2)}=(1/2,1/2,1,1).
\]
Predict, then update at each observation. The fine calculation gives
\[
p_1=\tfrac1{20}(15,3,2,0)^{\mathsf T},\qquad
p_2=\tfrac1{104}(41,15,40,8)^{\mathsf T}.
\]
\vfig{integrated_filter}
\captiontext{Within-packet probabilities change during the two observations, but the packet totals agree exactly with an independent two-state filter. The second coarse posterior is $(7/13,6/13)$.}
\head{Check the normalizers, not only the final vector}
The first and second predictive evidences are
\[
Z_1=\frac58,\qquad Z_2=\frac{13}{20},
\qquad Z_1Z_2=\frac{13}{32}.
\]
The unnormalized product $D_2KD_1Kp_0$ has mass $13/32$ and normalizes to $p_2$. Because the likelihood entries here lie in $[0,1]$, they can serve as probabilities of the specified observations. This makes the product their joint evidence in the stated model.

Two distinct equalities justify the coarse computation. First, $CK=K_cC$ transports prediction. Second, the likelihood is constant on each packet, so $CD_r=\overline D_rC$ transports observation weighting. Their composition preserves both the posterior packet totals and the evidence factors.

The coarse filter does not recover the within-packet shape. If a later observation distinguishes states inside a packet, those missing probabilities can become necessary, as the counterexample on \pagerefV{10} shows.
\takeaway{Exact filtering at a coarse level needs both dynamical lumpability and observation sufficiency. Agreement under one step does not excuse testing the other.}

\vpage{27}{Small at each level is not summable across levels}{REFINEMENT LIMITS / WHAT THE FINITE MODELS DO NOT PROVE}{Chapter VI, Part I: Corollary 10.5 and Section 17. Summability and uniform control are additional data.}
For each $n\ge1$, compare the two scalar operators $A_n=0$ and $B_n=1/n^2$ through the identity. Their ordinary defects have sizes $1/n^2$, while their square-root defects have sizes $1/n$:
\[
\sum_n|A_n-B_n|<\infty,
\qquad
\sum_n|A_n^{1/2}-B_n^{1/2}|=\infty.
\]
These are independent scalar comparisons, not a constructed inverse system. They isolate an error-scale obstruction without claiming a coherent refinement tower.
\vfig{summability}
\captiontext{Partial sums through $n=16$ illustrate the contrast. The convergence and divergence statements follow from the displayed series, not from the finite plot. A common positive gap gives the third, summable upper bound.}
\head{A uniform hypothesis can change the conclusion}
Replace the pair by $A_n=1$ and $B_n=1+1/n^2$. Then
\[
\left|B_n^{1/2}-A_n^{1/2}\right|
=\frac{1/n^2}{\sqrt{1+1/n^2}+1}
\le\frac1{2n^2}.
\]
This is the elementary form of the uniform-gap mechanism on \pagerefV{16}. Without such information, transferring a summable generator-defect budget through a fractional function can fail.

The formal release keeps three original research questions separate: realization of broader transfer semigroups, uniform stochastic-refinement control, and latent packet closure. The exact finite examples here do not settle their existence, domain, topology, or uniformity requirements.

Likewise, exact hidden-state elimination only yields the finite memory models that were supplied. It does not establish a finite latent realization for every memory law. A complete list of matrix identities can still leave an infinite-dimensional approximation problem untouched.
\takeaway{A finite example can expose a missing hypothesis. It cannot supply a projective-limit construction, a uniform constant, or a continuum realization merely by passing at several levels.}

\vpage{28}{How to use the evidence}{READER'S GUIDE / CALCULATION, DRAWING, AND FORMAL CONTRACT}{Chapter VI formal companion: thirty-three-result correspondence and reproduction record.}
The first 28 pages retain the worked mathematical atlas from the September release. Traditional Part I carries the corresponding core manuscript, including its figures and formal-status appendix; its theorem numbers are the references beneath these titles. The standalone guide uses internal visual-page links. The added pages beginning at Plate VI.29 address the strengthened branch in Traditional Part II.

The chosen release records all 33 scoped results complete, with 341 supporting theorem declarations in 47 proof modules and a saved audit of 782 namespace declarations. Those are inherited verification records, not a new Lean run and not counts of historically novel results. The three original conjectures retain their stated status.
\head{Read each kind of evidence at its actual strength}
\begin{tabularx}{\linewidth}{@{}>{\raggedright\arraybackslash}p{1.22in}YY@{}}
\toprule
Evidence & What is checked here & What it does not establish\\
\midrule
Exact matrices & Stochasticity, intertwining, projector identities, inverse formulas, and specific defects. & The same properties for arbitrary partitions or unbounded operators.\\
Rational algebra & Posterior/evidence identities, polynomial integrals, nodal errors, and complete finite path sums. & An unstated observation model or a continuum convergence theorem.\\
Symbolic formulas & The displayed exponentials, Laplace transforms, derivatives, and finite spectral powers under stated conditions. & A general-purpose formal proof of every analytic identity.\\
Vector drawings & Reproducible geometry and sampled curves from the explicit models. & Numerical quadrature certificates or an exhaustive infinite-time experiment.\\
Saved Lean records & The original contract, source files, dependency pins, and earlier verification logs. & Fresh kernel verification of this editorial addition.\\
\bottomrule
\end{tabularx}
\head{Reproduce a calculation before adapting it}
The visual-appendix tree contains the mathematical models, drawing generator, and exact-example verifier. Its README records their paths and execution commands. The inherited September verifier records 204 exact or symbolic assertions and 11 deliberately false controls rejected. The recorded monograph review reproduced that symbolic suite with the available SymPy and mpmath environment. Its recorded execution is identified on Plate VI.42. Its JSON record names every check and its finite scope.

The finite families are explicit: three-step Feynman--Kac paths are fully enumerated, selected rational priors are checked exhaustively within their stated denominator bound, and Green inverses are checked on four named grids. No finite count is used as a substitute for a general theorem.

To change a model, specify the state convention, partition, likelihood meaning, topology, and domain first. Then rerun the calculations and inspect the regenerated drawings. A preserved formula with changed assumptions is not automatically a preserved result.
\takeaway{Use the guide to see why a hypothesis matters and to compute a small case. Use the manuscript and its formal map for the scope of the general claim.}
